\subsection*{NT-E01. Cyclic reductions from precisely the division-field zeta functions}
Let $E/K$ be an elliptic curve over any number field, with or without
complex multiplication, and put $K_m=K(E[m])$ in a fixed algebraic closure.
Write $S_{E/K}(x)$ for the number of good prime ideals $\mathfrak p$ of
$K$ with $N\mathfrak p\le x$ for which $E_{\mathfrak p}(k_{\mathfrak p})$
is cyclic. The classical conditional density theorem gives
\[
 S_{E/K}(x)=\delta_{E/K}\frac{x}{\log x}
                 +o_{E,K}\!\left(\frac{x}{\log x}\right),\qquad
 \delta_{E/K}=\sum_{m\ge1}\frac{\mu(m)}{[K_m:K]}\ge0.       \tag{NT-E01a}
\]
The series is absolutely convergent. A sufficient analytic input is
\emph{RH for $\zeta_{K_m}$ for every squarefree $m$, including $m=1$}.
Thus $M$ supplies this consequence conditionally. Neither GRH for
$L(s,E)$ nor Artin holomorphy for the division-field Galois groups is an
additional hypothesis. No assertion of a positive density is implicit
in (NT-E01a).

The exact cited locators are Murty, \emph{On Artin's conjecture}
(1983), Theorem1, printed p.153, with its proof on pp.153--155;
and Campagna--Stevenhagen, \emph{Cyclic reduction densities for
elliptic curves} (2023), \S2, Lemma2.1, Corollary2.2 and
Propositions2.3--2.4 (author v2 pp.6--9). The latter verifies the
geometric degree growth, root-discriminant bound and large-prime
tail needed by Murty's theorem for arbitrary $E/K$.

\paragraph{Exact input identification.}
For a squarefree $m$, the compositum of the fields $K_\ell$ for
$\ell\mid m$ is $K_m$: the Chinese remainder decomposition of
$E[m]$ gives both inclusions. These are therefore exactly the fields
called $L_k$ in Murty's Theorem1. Its stated RH assumptions concern
their Dedekind zeta functions, rather than unspecified Artin factors.
The geometric criterion is that, at a good prime of characteristic
different from $\ell$, the full rank-two $\ell$-torsion is rational
over the residue field precisely when Frobenius acts trivially on
$E[\ell]$. Such full torsion prevents cyclicity. The characteristic
$p$ geometric $p$-torsion has rank at most one, so contributes no
additional rank-two obstruction. Corollary2.2 removes the resulting
finite set of residue-characteristic and model exceptions. This
identifies the inclusion--exclusion in (NT-E01a) with the stated
prime count. The limiting-density argument itself is the cited
classical theorem, not a new proof here.

Geometrically, each $K_m/K$ is the finite monodromy cover trivializing
the locally constant rank-two torsion sheaf of the elliptic scheme
over the appropriate arithmetic open. This uses finite quotients of
$\pi_1^{\text{ét}}$, and does not assert the Riemann hypothesis for
the infinite-image rank-two elliptic $L$-function. Work/run:
\texttt{GRH\_NIGHT/number\_theory}.

\subsection*{NT-E02. A finite Frobenius certificate and density-zero obstructions}
Suppose $E/K$ is non-CM. Let $T$ consist of the rational primes
dividing $30D_K$, those below a prime of bad reduction, and those for
which the image on $E[\ell]$ is not $\mathrm{GL}_2(\mathbb F_\ell)$.
This is a finite set, by the open image theorem. Put
$N=\prod_{\ell\in T}\ell$, $G_N=\operatorname{Gal}(K_N/K)$, and
\[
 C_N=\{g\in G_N:g|_{K_\ell}\ne1\text{ for every }\ell\mid N\}.
\]
Campagna--Stevenhagen Theorem1.1, author p.3, gives
\[
 \delta_{E/K}=\frac{|C_N|}{|G_N|}
   \prod_{\ell\nmid N}
       \left(1-\frac1{(\ell^2-1)(\ell^2-\ell)}\right).       \tag{NT-E02a}
\]
The product is strictly positive. Consequently $C_N=\varnothing$
is an unconditional finite certificate that all but finitely many
good reductions are noncyclic. If $C_N\ne\varnothing$, the hypothesis
of NT-E01 implies a positive density of cyclic reductions.

\paragraph{Complete finite-group deduction.}
For $m\mid N$, the event that a uniform $g\in G_N$ is trivial on
$K_m$ has probability $[K_m:K]^{-1}$. Intersections of these events
for primes dividing $m$ equal the $K_m$ event, by the compositum
identity proved above. Finite inclusion--exclusion therefore yields
\[
 \frac{|C_N|}{|G_N|}=\sum_{m\mid N}\frac{\mu(m)}{[K_m:K]}.
\]
The cited factorization gives (NT-E02a). For $\ell\ge2$, the
reciprocal of $(\ell^2-1)(\ell^2-\ell)$ is at most a constant
times $\ell^{-4}$ and is strictly less than one. The convergence
of its sum implies positivity of the product, by applying the
convergent logarithm series to all sufficiently large factors and
retaining the finitely many others.

If $C_N$ is empty, every element of $G_N$ lies in the kernel of at
least one restriction to $K_\ell$, $\ell\mid N$. Every good prime
outside the finite primes above $N$ and those ramified in $K_N$ has
such a Frobenius. Its reduction then contains
$(\mathbb Z/\ell\mathbb Z)^2$ for one such $\ell$, and is noncyclic.
This proves the finite-obstruction assertion without RH. Conversely,
$C_N\ne\varnothing$ makes the density expression positive, so NT-E01
gives infinitely many cyclic primes under its exact RH input.

For any $E/K$, including CM curves, the simpler obstruction
$E[\ell]\subset E(K)$ prevents cyclic reduction at every good prime
away from $\ell$, by the same reduction argument. Full rational
two-torsion is thus an immediate obstruction. Over a general $K$
the absence of this particular obstruction is not a positivity
criterion: the finite kernels in (NT-E02a) may jointly cover $G_N$
without any one being all of $G_N$. Actual elliptic examples of this
phenomenon are supplied by Campagna--Stevenhagen, Theorem5.1 and
\S5 (author pp.14--16); their construction uses three quadratic
division fields with Klein-four compositum. No independence of
division fields is silently assumed.

This is a finite monodromy-cover certificate, not a criterion for
a prescribed rational point to generate every cyclic reduction.
It also makes no claim about prime values of $\#E(\mathbb F_p)$.
Work/run: \texttt{GRH\_NIGHT/number\_theory}.

\subsection*{NT-E03. Lang--Trotter upper bounds without an extra holomorphy conjecture}
For a fixed non-CM elliptic curve $E/\mathbb Q$, set
$a_p=p+1-\#E(\mathbb F_p)$ at good primes, and
$P_{E,a}(x)=\#\{p\le x:p\nmid N_E,\ a_p=a\}$.
Zywina, \emph{Bounds for the Lang--Trotter conjectures},
Theorem1.2, author p.1, proves
\[
 P_{E,a}(x)\ll_E\frac{x^{4/5}}{(\log x)^{3/5}},\qquad
 P_{E,0}(x)\ll_E\frac{x^{3/4}}{(\log x)^{1/2}}.            \tag{NT-E03a}
\]
The first constant is uniform over integers $a$. These are upper
bounds, not the Lang--Trotter asymptotic. The source defines its GRH
on author p.3 as Dedekind RH for number fields. The proof permits the
following sufficient smaller input: adopted $A$, together with RH
for $\zeta_{\mathbb Q(E[\ell])}$ for every sufficiently large prime
$\ell$. Hence $M$ supplies the conclusion. Non-CM $\mathrm{GL}_2$
GRH is not being inferred.

\paragraph{Complete holomorphy/input audit.}
For large $\ell$ the mod-$\ell$ image is $G=\mathrm{GL}_2(\mathbb F_\ell)$.
Write $L=\mathbb Q(E[\ell])$, let $B$ be its upper triangular
subgroup, and let $U$ be its unipotent subgroup. Zywina's complete
\S4 proof, author pp.10--11, applies its Theorem2.3 to
$L^U/L^B$. Its group $B/U$ is abelian, so its Artin functions are
Hecke functions and the holomorphy assumption of Theorem2.3 is
known. Moreover $L/L^U$ has the cyclic group $U$, so the all-prime
abelian factorization NT-H11 gives
$\zeta_L=\zeta_{L^U}\prod_{\chi\ne1}L_{L^U}(s,\chi)$ with every
other factor entire. Thus RH for $\zeta_L$ implies RH for
$\zeta_{L^U}$, exactly what Theorem2.3 requires.

For $a=0$ the proof replaces $U$ by the subgroup $H$ of triangular
matrices with equal diagonal entries. This group is abelian:
writing a matrix as $aI+b\left(\begin{smallmatrix}0&1\\0&0\end{smallmatrix}\right)$
makes commutativity explicit. The same factorization transfers RH
from $\zeta_L$ to $\zeta_{L^H}$, while $B/H$ is again abelian.
The separate step choosing a prime $\ell$ split in
$\mathbb Q(\sqrt{a_p^2-4p})$ uses only RH for quadratic Dedekind
zeta functions and for $\zeta_{\mathbb Q}$: see Murty--Murty--Saradha
(1988), Lemma4.4 and its entire proof, printed pp.269--270.
These are supplied by $A$. This proves the stated sufficient input
reduction. The quantitative estimates in (NT-E03a) remain the
cited classical theorem.

The geometric data used here are finite mod-$\ell$ monodromy of
$H^1_{\text{ét}}(E_{\overline{\mathbb Q}},\mathbb F_\ell)$, with
the usual dual convention immaterial to the division field, and
abelian covers of intermediate arithmetic bases. An arbitrary
nonabelian use of Zywina's Theorem2.3 would still require its stated
Artin holomorphy hypothesis. Work/run:
\texttt{GRH\_NIGHT/number\_theory}.

\subsection*{NT-E04. Uniform counts with a prescribed Frobenius field}
Let $E/\mathbb Q$ be non-CM and let $k$ be an imaginary quadratic
field with class number $h_k$ and absolute discriminant $d_k$.
Put
\[
 P_{E,k}(x)=\#\{p\le x:p\nmid N_E,\
            \mathbb Q(\sqrt{a_p^2-4p})\cong k\}.
\]
Zywina Theorem1.3(i), author p.2, gives under Dedekind GRH
\[
 P_{E,k}(x)\ll_E
   h_k^{-3/5}\frac{x^{4/5}}{(\log x)^{3/5}}
                   +x^{1/2}(\log x)^3,                    \tag{NT-E04a}
\]
with a constant independent of $k$. Its proof is \S5.1,
author pp.13--14, after the algebraic reductions on pp.11--13.
Lemma5.5, author p.13, also gives $P_{E,k}(x)=0$ when $d_k>4x$.

For precision, this field family is larger than the division
fields in NT-E03. Let $H_k$ be the Hilbert class field of $k$,
and for split primes $\ell\ge5$ let $L_{k,\ell}/H_k$ be the
field of the joint representation formed from $E[\ell]$ and
the modulus-$\ell$ ray-class representation
$\psi_{k,\ell}:G_{H_k}\to
(\mathcal O_k/\ell\mathcal O_k)^*/\mathcal O_k^*$.
Zywina \S3.2--3.3, pp.7--10, defines it. The proof uses RH for
the Dedekind zeta functions of the auxiliary fields
$L_{k,\ell}^{\mathcal U}$ appearing in \S5, together with RH
for $\zeta_k$ and $\zeta_{\mathbb Q}$ in the split-prime choice.
Here $\mathcal B$ is the preimage of the triangular subgroup and
$\mathcal U$ is the image of the pairs $(aI+bN,a)$; the quotient
$\mathcal B/\mathcal U$ is abelian, so holomorphy is known.
This is an exact sufficient family specified by the cited proof;
$M$ supplies all of it. We do not replace it by Dirichlet GRH alone.

These counts distinguish endomorphism fields of reductions using
finite torsion and ray-class monodromy. They do not establish the
conjectured $x^{1/2}/\log x$ asymptotic. Work/run:
\texttt{GRH\_NIGHT/number\_theory}.

\subsection*{NT-E05. A refined lower bound for distinct Frobenius fields}
Let $\mathcal D_E(x)$ be the set of imaginary quadratic fields
arising as $\mathbb Q(\sqrt{a_p^2-4p})$ at good primes $p\le x$
for a fixed non-CM $E/\mathbb Q$. Under the exact input of NT-E04
for every imaginary quadratic $k$, together with adopted $A$,
\[
 |\mathcal D_E(x)|\gg_E
 \frac{x^{2/7}}
      {(\log x)^{4/7}(\log\log x)^{6/7}}
 \qquad(x\longrightarrow\infty).                         \tag{NT-E05a}
\]
In particular $M$ supplies these inputs. This is a newly proved
formulation in this run, not a claim of literature novelty.
Zywina's Corollary1.4, author p.2 with proof p.15, gives
$x^{2/7}/(\log x)^{10/7}$ using a weaker class-number estimate;
the proof below retains the sharper checked NT-H02 input.

\paragraph{Complete proof.}
Put $\Lambda(t)=\log\log(e^e t)$ for $t\ge1$. The explicit
NT-H02 bound, together with its finitely many excluded fundamental
discriminants, gives an absolute $c>0$ such that
\[
 h_k\ge c\,\frac{\sqrt{d_k}}{\Lambda(d_k)}
 \quad\hbox{for every imaginary quadratic }k.             \tag{NT-E05b}
\]
Indeed the displayed $V(d)$ in NT-H01 is $O(\log\log d)$ as
$d\to\infty$; lowering $c$ handles the finitely many smaller
discriminants, all of whose class numbers are positive.

Let $D=|\mathcal D_E(x)|$ and take $x$ sufficiently large that
$D\ge1$. Every such field has $d_k\le4x$ by NT-E04. Distinct
imaginary quadratic fields have distinct positive absolute
fundamental discriminants, so, ordering these discriminants,
the $j$th is at least $j$. It follows that
\begin{align*}
 \sum_{k\in\mathcal D_E(x)}h_k^{-3/5}
 &\ll \Lambda(4x)^{3/5}
                \sum_{k\in\mathcal D_E(x)}d_k^{-3/10}\\
 &\ll \Lambda(x)^{3/5}\sum_{j=1}^D j^{-3/10}
 \ll \Lambda(x)^{3/5}D^{7/10}.
\end{align*}
The last estimate uses
$\sum_{j=1}^D j^{-3/10}\le(10/7)D^{7/10}$, directly by integration of
$t^{-3/10}$ on $[1,D]$; its implicit constant is independent of
$E$, $x$ and the chosen fields.

Every good prime contributes to exactly one of these fields.
The prime number theorem, removal of the finitely many primes
dividing $N_E$, and (NT-E04a) therefore give
\[
 \frac{x}{\log x}\ll_E
    D^{7/10}\frac{x^{4/5}\Lambda(x)^{3/5}}{(\log x)^{3/5}}
                      +D x^{1/2}(\log x)^3.               \tag{NT-E05c}
\]
At least one right-hand term is at least a fixed positive
multiple of $x/\log x$. If it is the first, solving the
inequality gives
\[
 D\gg_E x^{2/7}(\log x)^{-4/7}\Lambda(x)^{-6/7}.
\]
If it is the second, then $D\gg_E x^{1/2}/(\log x)^4$.
The ratio of this latter lower bound to the former tends to
infinity, since it is
\[
 x^{3/14}(\log x)^{-24/7}\Lambda(x)^{6/7}\longrightarrow\infty.
\]
Thus the first lower bound follows in both cases for all
sufficiently large $x$. Finally
$\Lambda(x)\sim\log\log x$, which proves (NT-E05a).

This measures the diversity of Frobenius endomorphism fields
among closed fibres of an elliptic scheme. It does not prove
prime-order reductions or the Koblitz conjecture. Work/run:
\texttt{GRH\_NIGHT/number\_theory}.
